\section{The space of proofs of one lemma}

\lean{RequestProject/Variations.lean} proves each result of the development several
times, moving along the axes named in the proof-writing guides.  Those axes span a
Boolean lattice, so the catalogue of variants is literally a set of vertices of a cube ---
and one can see at a glance which corners are occupied.

\begin{figure}[ht]
\centering
\resizebox{\textwidth}{!}{%
\begin{tikzpicture}[x=1cm,y=1cm,font=\scriptsize]
  \tikzset{v/.style={circle,draw=black!65,fill=white,inner sep=1.6pt},
           occ/.style={v,fill=tagglue!45,draw=tagglue},
           emp/.style={v,fill=white,draw=black!35,densely dotted}}
  \def\dx{7.0}\def\dy{4.4}\def\dz{3.0}
  \coordinate (o000) at (0,0);
  \coordinate (o100) at (\dx,0);
  \coordinate (o010) at (0,\dy);
  \coordinate (o110) at (\dx,\dy);
  \coordinate (o001) at (\dz,{.62*\dz});
  \coordinate (o101) at ({\dx+\dz},{.62*\dz});
  \coordinate (o011) at (\dz,{\dy+.62*\dz});
  \coordinate (o111) at ({\dx+\dz},{\dy+.62*\dz});
  \foreach \a/\b in {o000/o100,o000/o010,o100/o110,o010/o110,
                     o001/o101,o001/o011,o101/o111,o011/o111,
                     o000/o001,o100/o101,o010/o011,o110/o111}
    {\draw[black!45] (\a) -- (\b);}
  \node[occ] at (o000) {}; \node[occ] at (o100) {};
  \node[occ] at (o010) {}; \node[occ] at (o110) {};
  \node[occ] at (o001) {}; \node[occ] at (o101) {};
  \node[emp] at (o011) {}; \node[occ] at (o111) {};
  \node[anchor=north east,align=right] at ($(o000)+(-.15,-.15)$)
    {\textbf{term, manual, one-liner}\\ \lean{cofinal\_v₁}, \lean{diagonal\_v₂}};
  \node[anchor=north west,align=left] at ($(o100)+(.15,-.15)$)
    {\textbf{tactic, manual, one-liner}\\ \lean{cofinal\_v₆} (\lean{<;>}),
     \lean{diagonal\_v₃}};
  \node[anchor=south east,align=right] at ($(o010)+(-.15,.15)$)
    {\textbf{term, manual, expanded}\\ \lean{cofinal\_v₂}};
  \node[anchor=south east,align=right] at ($(o110)+(-.15,.2)$)
    {\textbf{tactic, manual, expanded}\\ \lean{cofinal\_v₃},
     \lean{cofinal\_v₅} (\lean{calc})};
  \node[anchor=east,align=right,text=tagstructural] (labA) at (-0.5,2.6)
    {\textbf{term, automated}\\ \lean{ennreal\_mul\_pointfree},\\
     \lean{ennreal\_mul\_swap} (\lean{Function.swap})};
  \draw[thin,tagstructural!60] (labA.east) -- (o001);
  \node[anchor=west,align=left,text=tagstructural] at ($(o101)+(.25,-.1)$)
    {\textbf{tactic, automated, one-liner}\\ \lean{cofinal\_v₄} (\lean{simp only}),\\
     \lean{diagonal\_v₄} (\lean{rw}), \lean{grw}};
  \node[anchor=south east,align=right,text=black!45] at ($(o011)+(-.25,.25)$)
    {\emph{unoccupied}: automated\\ \emph{and} verbose in term mode};
  \node[anchor=west,align=left] at ($(o111)+(.25,.15)$)
    {\textbf{tactic, automated, expanded}\\ \lean{cardinal\_v₃} (\lean{simpa using})};
  % axes legend
  \begin{scope}[shift={(-1.2,-3.4)}]
    \draw[-{Stealth[length=2mm]},black!60] (0,0) -- ++(1.4,0)
      node[right,font=\scriptsize] {style: term $\to$ tactic};
    \draw[-{Stealth[length=2mm]},black!60] (0,0) -- ++(0,1.0)
      node[above,font=\scriptsize] {granularity};
    \draw[-{Stealth[length=2mm]},black!60] (0,0) -- ++(.85,.53)
      node[right,font=\scriptsize,xshift=1pt] {automation};
  \end{scope}
\end{tikzpicture}}
\caption{The design space of the mutual-cofinality lemma as the Boolean lattice
$B_3$.  Green vertices are realised by compiled proofs in \lean{Variations.lean}.}
\end{figure}

\subsection{A fourth axis: duality}

Duality is not a point of the cube but an involution acting on the whole of it: applying
$(-)^{\mathrm{od}}$ to any proof yields a proof of the $\Inf$-statement, and Lean accepts
it with a single type ascription.

\begin{figure}[ht]
\centering
\begin{tikzpicture}[x=1cm,y=1cm,font=\scriptsize]
  \draw[rounded corners=8pt,fill=tagglue!7,draw=tagglue!50] (-4.1,-1.4) rectangle (0.1,1.4);
  \draw[rounded corners=8pt,fill=tagreducible!7,draw=tagreducible!50]
    (4.1,-1.4) rectangle (8.3,1.4);
  \node[align=center] at (-2.0,.75) {$\Sup$-world\\ \lean{iSup₂\_eq\_iSup\_diagonal}};
  \node[align=center] at (-2.0,-.55) {\lean{iSup\_op\_iSup\_diagonal}\\
    \lean{ciSup₂\_eq\_ciSup\_diagonal}};
  \node[align=center] at (6.2,.75) {$\Inf$-world\\ \lean{iInf₂\_eq\_iInf\_diagonal}};
  \node[align=center] at (6.2,-.55) {\lean{iInf\_op\_iInf\_diagonal}\\
    \lean{ciInf₂\_eq\_ciInf\_diagonal}};
  \draw[{Stealth[length=2.4mm]}-{Stealth[length=2.4mm]},bridgecol,line width=1pt]
    (0.35,0) -- (3.85,0);
  \node[above=2pt,text=bridgecol,font=\scriptsize] at (2.1,0) {\lean{(α := αᵒᵈ)}};
  \node[below=2pt,text=bridgecol,font=\scriptsize] at (2.1,0) {cost: 0 lines};
\end{tikzpicture}
\caption{The dual half of \lean{Diagonal.lean} is generated, not written.}
\end{figure}

\subsection{From the abstraction to the applications}

\begin{figure}[ht]
\centering
\begin{tikzpicture}[x=1cm,y=1cm,font=\scriptsize,
  every node/.style={inner sep=2.5pt}]
  \tikzset{core/.style={rounded corners=3pt,draw=tagstructural,fill=tagstructural!10,
      font=\ttfamily\scriptsize,align=center},
    appl/.style={rounded corners=3pt,draw=tagglue,fill=tagglue!10,
      font=\ttfamily\scriptsize,align=center},
    inp/.style={draw=black!35,fill=white,font=\scriptsize\itshape,align=center,
      rounded corners=2pt}}
  \node[core] (op)  at (0,2.0) {iSup\_op\_iSup\_diagonal};
  \node[core] (cop) at (0,-1.6) {ciSup\_op\_ciSup\_diagonal};
  \node[inp] (d1) at (-4.3,3.3) {ENNReal.iSup\_add\\ ENNReal.add\_iSup};
  \node[inp] (d2) at (-4.3,1.9) {ENNReal.iSup\_mul\\ ENNReal.mul\_iSup};
  \node[inp] (d3) at (-4.3,0.6) {ENat.iSup\_add\\ ENat.add\_iSup};
  \node[inp] (d4) at (-4.3,-1.6) {Cardinal.ciSup\_add\_ciSup\\ + bddAbove\_range\_add};
  \node[appl] (a1) at (5.0,3.3) {ennreal\_iSup\_add\_iSup};
  \node[appl] (a2) at (5.0,2.0) {ennreal\_iSup\_mul\_iSup};
  \node[appl] (a3) at (5.0,0.7) {enat\_iSup\_add\_iSup};
  \node[appl] (a4) at (5.0,-1.6) {cardinal\_ciSup\_add\_\\ciSup\_diagonal};
  \node[appl] (a5) at (5.0,-3.2) {cardinal\_ciSup\_add\_\\ciSup\_of\_monotone};
  \node[appl] (a6) at (0.2,4.5) {ennreal\_tendsto\_add\_atTop};
  \node[core] (bim) at (0,3.5) {..\_of\_bimonotone};
  \foreach \s in {d1,d2,d3}{\draw[use] (\s) -- (op);}
  \draw[use] (d4) -- (cop);
  \foreach \t in {a1,a2,a3}{\draw[use] (op) -- (\t);}
  \draw[use] (cop) -- (a4);
  \draw[use] (a4) -- (a5);
  \draw[use] (bim) -- (a6);
  \node[font=\scriptsize,text=black!55,align=center] at (-4.3,-3.0)
    {\emph{inputs}: the distributivity\\ laws of each theory};
  \node[font=\scriptsize,text=black!55] at (5.0,-4.2) {\emph{outputs}};
\end{tikzpicture}
\caption{The operation-level packaging turns each application into ``supply the two
distributivity lemmas''.  The multiplicative statement for $\enn$ and the additive one
for $\mathbb{N}_\infty$ cost two lines each, once the abstraction exists.}
\end{figure}

\subsection{Inventory}

\begin{center}\footnotesize
\begin{tabular}{@{}llr@{}}
\toprule
Statement & Axes exercised & Variants \\
\midrule
mutual cofinality & term/tactic, forward/backward, \lean{calc}, \lean{<;>} & 7 \\
diagonal corollary & corollary, standalone, \lean{rw}, duality & 6 \\
binary-operation form & \lean{trans} chain, \lean{simp\_rw}, point-free & 4 \\
boundedness helpers & set-level vs.\ family-level & 5 \\
conditionally complete collapse & currying-first vs.\ \lean{le\_antisymm}-first & 5 \\
\lean{Cardinal} payoff & \lean{rw}, \lean{simpa using}, linear-order bridge & 5 \\
$\enn$ & operation-level, categorical, \lean{sSupHom} transport & 6 \\
structural atoms & term/tactic, manual/automated & 6 \\
\bottomrule
\end{tabular}
\end{center}

\clearpage
